\documentclass[letterpaper, 12pt]{article}
% ============================================================
%  CONFIGURACION COMUN PARA SOLUCIONARIOS PEP 1
% ============================================================

\usepackage[utf8]{inputenc}
\usepackage[spanish]{babel}
\usepackage{amsmath, amssymb, amsfonts}
\usepackage{geometry}
\usepackage{fancyhdr}
\usepackage{enumitem}
\usepackage{graphicx}
\usepackage{hyperref}
\usepackage{needspace}
\usepackage{parskip}
\usepackage{xcolor}
\usepackage{tcolorbox}
\usepackage{eso-pic}
\usepackage{tikz}
\tcbuselibrary{skins, breakable}

\geometry{letterpaper, margin=1.7cm, top=3cm, bottom=2.5cm, headheight=2cm}

\definecolor{celesteSuave}{HTML}{F0F7FF}
\definecolor{azulFuerte}{HTML}{1A5276}
\definecolor{enlace}{HTML}{1A5276}

\newcommand{\R}{\mathbb{R}}
\newcommand{\Dom}{\operatorname{Dom}}
\newcommand{\Rec}{\operatorname{Rec}}
\newcommand{\Cod}{\operatorname{Cod}}

\newcommand{\logoup}{../../../../../template/images/logo_up.png}
\newcommand{\logousach}{../../../../../template/images/logo_usach.png}
\newcommand{\logofondo}{../../../../../template/images/logo_up.png}
\newcommand{\institucion}{Usach Premium}
\newcommand{\curso}{C\'alculo I}
\newcommand{\autor}{Jaime Riquelme}
\newcommand{\semestre}{1er. Semestre de 2026}
\newcommand{\fraseportada}{``Por y para estudiantes''}
\newcommand{\webinstitucion}{www.usachpremium.cl}
\newcommand{\instagraminstitucion}{https://www.instagram.com/usach.premium}
\newcommand{\transparencia}{0.05}

\hypersetup{
    colorlinks=true,
    linkcolor=enlace,
    urlcolor=enlace,
    citecolor=enlace,
    pdfborder={0 0 0},
}

\pagestyle{fancy}
\fancyhf{}
\fancyhead[L]{\includegraphics[height=1.2cm]{\logoup}}
\fancyhead[R]{%
    \begin{minipage}[b]{0.42\textwidth}
        \raggedleft
        \fontsize{9pt}{11pt}\selectfont
        \textbf{\color{azulFuerte}\institucion} \\
        \curso\ -- Solucionario PEP 1 \\
        \textit{\color{gray}@usach.premium}
    \end{minipage}\hspace{0.1cm}%
    \includegraphics[height=1.2cm]{\logousach}%
}
\fancyfoot[L]{\fontsize{9pt}{11pt}\selectfont \textit{\fraseportada}}
\fancyfoot[R]{\fontsize{9pt}{11pt}\selectfont P\'agina \thepage}
\renewcommand{\headrulewidth}{0.4pt}
\renewcommand{\footrulewidth}{0.4pt}

\fancypagestyle{plain}{
    \fancyhf{}
    \renewcommand{\headrulewidth}{0pt}
    \renewcommand{\footrulewidth}{0pt}
}

\newcommand{\MarcaAgua}{
    \AddToShipoutPicture{
        \AtPageCenter{
            \begin{tikzpicture}[remember picture, overlay]
                \node[opacity=\transparencia] at (0,0)
                    {\includegraphics[width=0.7\textwidth]{\logofondo}};
            \end{tikzpicture}
        }
    }
}

\newtcolorbox{solucionbox}[1]{
    colback=white,
    colframe=azulFuerte,
    coltitle=white,
    colbacktitle=azulFuerte,
    title=\textbf{#1},
    title after break=\textbf{#1} (continuaci\'on),
    fonttitle=\bfseries,
    arc=4pt,
    boxrule=1pt,
    enhanced,
    breakable,
    before={\par\Needspace{0.36\textheight}}
}

\newtcolorbox{respuestabox}{
    colback=celesteSuave,
    colframe=azulFuerte,
    boxrule=0.8pt,
    arc=4pt
}

\newcommand{\portadasolucionario}[2]{
    \thispagestyle{plain}
    \AddToShipoutPicture*{%
        \put(54, 700){\includegraphics[width=2cm]{\logoup}}
        \put(420, 706){\includegraphics[width=5cm]{\logousach}}
    }
    \vspace*{0.5cm}
    \begin{center}
        \color{azulFuerte}
        {\huge \textbf{SOLUCIONARIO}} \\[0.5cm]
        {\Huge \textbf{\curso}} \\[0.3cm]
        {\Large #1} \\[0.3cm]
        {\large #2} \\[1.4cm]
        \begin{tcolorbox}[colback=celesteSuave, colframe=azulFuerte, arc=10pt,
            width=0.82\textwidth, center, boxrule=1.5pt]
            \centering
            \Large\textbf{Desarrollo paso a paso}
        \end{tcolorbox}
    \end{center}
    \vspace{1.4cm}
    \begin{center}
        {\Large \textbf{\semestre}} \\[0.7cm]
        {\large \textbf{\autor\ -- \institucion}}
    \end{center}
    \vfill
    \begin{center}
        {\color{azulFuerte}\hrule height 0.5pt}
        \vspace{0.3cm}
        \begin{minipage}{0.45\textwidth}
            \centering
            \textbf{Instagram:} \\
            \small\href{\instagraminstitucion}{@usach.premium}
        \end{minipage}
        \begin{minipage}{0.45\textwidth}
            \centering
            \textbf{Web:} \\
            \small\href{http://\webinstitucion}{\webinstitucion}
        \end{minipage}
    \end{center}
    \newpage
    \MarcaAgua
}

\newcommand{\respuesta}[1]{
    \begin{respuestabox}
        \textbf{Respuesta:} #1
    \end{respuestabox}
}


\begin{document}
\portadasolucionario{Gu\'ia de Preparaci\'on PEP 1}{Secci\'on IV -- Par\'ametros e intersecciones}

\begin{solucionbox}{Ejercicio 4.1}
Determine todos los valores de $k\in\R$ para que la par\'abola
\[
    y=kx^2+4x+1
\]
y la recta
\[
    y=(k+2)x-3
\]
se intersecten en dos puntos diferentes.

\textbf{Soluci\'on:}

Para encontrar los puntos de intersecci\'on, igualamos ambas expresiones:
\[
    kx^2+4x+1=(k+2)x-3.
\]
Llevamos todo a un lado:
\[
    kx^2+4x-(k+2)x+1+3=0.
\]
Reduciendo t\'erminos semejantes:
\[
    kx^2+(2-k)x+4=0.
\]

Las intersecciones entre la par\'abola y la recta corresponden a las soluciones de esta ecuaci\'on. Para que haya dos puntos diferentes, necesitamos que la ecuaci\'on sea cuadr\'atica y tenga dos soluciones reales distintas.

Primero, debe cumplirse
\[
    k\neq 0.
\]
Luego, su discriminante debe ser positivo:
\[
    \Delta>0.
\]
Calculamos:
\[
    \Delta=(2-k)^2-4(k)(4).
\]
Entonces
\[
    \Delta=(2-k)^2-16k.
\]
Desarrollando:
\[
    \Delta=k^2-4k+4-16k
    =k^2-20k+4.
\]
Por lo tanto, debemos resolver
\[
    k^2-20k+4>0.
\]
Resolvemos la ecuaci\'on asociada:
\[
    k^2-20k+4=0.
\]
Por f\'ormula cuadr\'atica:
\[
    k=\frac{20\pm\sqrt{(-20)^2-4(1)(4)}}{2}
     =\frac{20\pm\sqrt{400-16}}{2}.
\]
Luego,
\[
    k=\frac{20\pm\sqrt{384}}{2}
     =\frac{20\pm 8\sqrt6}{2}
     =10\pm4\sqrt6.
\]
Como la par\'abola en $k$ abre hacia arriba,
\[
    k^2-20k+4>0
\]
cuando
\[
    k<10-4\sqrt6
    \quad\vee\quad
    k>10+4\sqrt6.
\]
Finalmente, recordando que $k\neq0$, se obtiene
\[
    k\in(-\infty,0)\cup(0,10-4\sqrt6)\cup(10+4\sqrt6,+\infty).
\]

\respuesta{$k\in(-\infty,0)\cup(0,10-4\sqrt6)\cup(10+4\sqrt6,+\infty)$.}
\end{solucionbox}

\begin{solucionbox}{Ejercicio 4.2}
Considere las par\'abolas
\[
    C_1:\ y=(k+1)x^2-3x+k,
    \qquad
    C_2:\ y=2x^2+(k-2)x+1.
\]
Determine los valores de $k\in\R$ para que $C_1$ y $C_2$ se intersecten en dos puntos diferentes.

\textbf{Soluci\'on:}

Para hallar las intersecciones, igualamos ambas ecuaciones:
\[
    (k+1)x^2-3x+k=2x^2+(k-2)x+1.
\]
Llevamos todo a un lado:
\[
    (k+1)x^2-2x^2-3x-(k-2)x+k-1=0.
\]
Reducimos:
\[
    (k-1)x^2+(-3-k+2)x+(k-1)=0.
\]
Es decir,
\[
    (k-1)x^2-(k+1)x+(k-1)=0.
\]

Para que las par\'abolas se intersecten en dos puntos diferentes, la ecuaci\'on debe ser cuadr\'atica y tener dos soluciones reales distintas.

Primero:
\[
    k-1\neq0
    \quad\Longrightarrow\quad
    k\neq1.
\]
Luego, exigimos
\[
    \Delta>0.
\]
Calculamos:
\[
    \Delta=(-(k+1))^2-4(k-1)(k-1).
\]
Por lo tanto,
\[
    \Delta=(k+1)^2-4(k-1)^2.
\]
Desarrollando:
\[
    \Delta=(k^2+2k+1)-4(k^2-2k+1).
\]
Entonces
\[
    \Delta=k^2+2k+1-4k^2+8k-4
    =-3k^2+10k-3.
\]
Debemos resolver
\[
    -3k^2+10k-3>0.
\]
Multiplicando por $-1$ y cambiando el sentido:
\[
    3k^2-10k+3<0.
\]
Factorizamos:
\[
    3k^2-10k+3=(3k-1)(k-3).
\]
Luego,
\[
    (3k-1)(k-3)<0.
\]
La expresi\'on es negativa entre sus ra\'ices:
\[
    \frac13<k<3.
\]
Finalmente, debemos excluir $k=1$:
\[
    k\in\left(\frac13,1\right)\cup(1,3).
\]

\respuesta{$k\in\left(\dfrac13,1\right)\cup(1,3)$.}
\end{solucionbox}

\end{document}
